% ===================================================================
%  CADRO N.º 6 — A casa por dentro, o mobiliario, a alimentación,
%  a iluminación.
%  Traduction 2026 d'apres texto/io, controlee sur texto/fr, texto/es
%  et texto/pt. Consigne : voir texto/gl/10-cadro-01.tex.
%
%  LE PLUS LONG TABLEAU DE LA COLONNE, et celui qui nomme le plus
%  d'objets domestiques : cinq pieces, deux cent quarante renvois.
%  C'est donc le second banc d'essai du tableau 2, et il donne le
%  meme resultat. Ce qui se manie tous les jours a son mot galicien :
%  « vasoira » le balai, « billa » le robinet, « tixola » la poele,
%  « fregadeiro » l'evier, « colcha » la courtepointe, « sabas » les
%  draps, « peite » le peigne, « mistos » les allumettes.
%
%  ET LE MOT LE PLUS UTILE DE CE TABLEAU N'EST DANS AUCUNE DES DEUX
%  VOISINES SOUS CETTE FORME : « mistos » pour les allumettes, la ou
%  l'espagnol dit « cerillas » et le portugais « fósforos ». Trois
%  langues, trois mots, pour l'objet le plus banal de la cuisine.
%
%  L'ECART DECLARE DU BLOC c1-06-1. Le fac-simile ido y ouvre sur
%  « (6) », qui est le savon de l'alinea 2 ; c'est la femme de chambre
%  « (16) », comme le porte le francais et comme la planche le grave.
%  La colonne le rend donc, avec les dix qui la precedent, et
%  renvoji.py le passe sans rien dire : c'est le premier de ses trois
%  ecarts declares.
% ===================================================================

%%K t06-apar-1 apar
\VUsaut{6.0mm}
\VUtitre{15.0pt}{60}{CADRO N.\textsuperscript{o} 6}
\VUsaut{1.4mm}
\VUtitre{11.4pt}{0}{A casa por dentro, o mobiliario, a alimentación,}
\VUsaut{0.4mm}
\VUtitre{11.4pt}{0}{a iluminación.}
\VUsaut{2.2mm}

%%K t06-apar-2 apar
A casa está rematada. O tío de Xoán está instalado na súa nova vivenda, da que xa
coñecemos a distribución. Vexamos agora como amoblou a súa casa. Visitaremos
primeiro:

%%K t06-tit-1 sub
\VUpk{10.2pt}{40}{O dormitorio.}
\VUsaut{3.0mm}

%%K t06-c1-01-1 p
1. --- Chegamos a mal tempo, pois a doncela está ocupada limpando o \VUgras{cuarto.}

%%K t06-c1-02-1 p
2. --- Sobre o \VUgras{tocador} \textsuperscript{(1)} hai aquí e alá distintos
obxectos cos que un se lava, se peitea, en poucas palabras se arranxa. Son a
\VUgras{palangana} \textsuperscript{(2)} e a \VUgras{xerra da auga}
\textsuperscript{(3)}, o \VUgras{cepillo do pelo} \textsuperscript{(4)}, o
\VUgras{peite} \textsuperscript{(5)}, o \VUgras{xabón} \textsuperscript{(6)} e a
\VUgras{esponxa} \textsuperscript{(7)}, e finalmente un \VUgras{frasco} limpo
\textsuperscript{(8)} para o dentífrico líquido.

%%K t06-c1-03-1 p
3. --- No chan, diante do tocador, velaquí o \VUgras{balde} \textsuperscript{(9)}
para recoller as augas sucias.

%%K t06-c1-04-1 p
4. --- Vemos na \VUgras{antecámara} \textsuperscript{(10)} algúns
\VUgras{colgadoiros} \textsuperscript{(13)} e dous \VUgras{paraugas}
\textsuperscript{(11)} no \VUgras{paragüeiro} \textsuperscript{(12)}.

%%K t06-c1-05-1 p
5. --- Do outro lado da porta está o \VUgras{armario de espello}
\textsuperscript{(14)} que garda a \VUgras{roupa branca} \textsuperscript{(15)}.

%%K t06-c1-06-1 p
6. --- Aproveitemos que a \VUgras{doncela} \textsuperscript{(16)} está a facer a
\VUgras{cama} \textsuperscript{(17)} para examinar as súas distintas partes. O
\VUgras{armazón da cama} \textsuperscript{(20)} leva un bo \VUgras{somier de
resortes} \textsuperscript{(21)} sobre o que Xustina vai pór axiña un
\VUgras{colchón} de la \textsuperscript{(22)}. Despois collerá das cadeiras as dúas
\VUgras{sabas} \textsuperscript{(23)} e a \VUgras{manta} \textsuperscript{(24)}.
Entón porá na cabeceira o \VUgras{traveseiro} \textsuperscript{(25)} e a
\VUgras{almofada} \textsuperscript{(26)}, que están co \VUgras{edredón}
\textsuperscript{(27)} sobre a \VUgras{colcha} \textsuperscript{(32)}. Usa un
plumeiro para quitar o po ás \VUgras{cortinas} \textsuperscript{(19)} e ao
\VUgras{doselete} \textsuperscript{(18)}, e velaquí unha parte do traballo rematada.

%%K t06-c1-07-1 p
7. --- Despois terá que arranxar un pouco a \VUgras{mesiña de noite}
\textsuperscript{(28)}, dar corda ao \VUgras{espertador} \textsuperscript{(29)},
preparar a \VUgras{lampadiña de noite} \textsuperscript{(30)} e sacar a
\VUgras{vela} \textsuperscript{(31)} para limpar a palmatoria.

%%K t06-tit-2 sub
\VUsaut{5.0mm}
\VUpk{10.2pt}{40}{O cesto da costura e os apeiros da lareira.}
\VUsaut{3.0mm}

%%K t06-c1-08-1 p
8. --- O \VUgras{cesto da costura} \textsuperscript{(34)} que está preto do
\VUgras{tallo} \textsuperscript{(33)} tamén precisa moito de ser arranxado. Terá que
dobrar os \VUgras{bordados} \textsuperscript{(35)}, gardar na \VUgras{mesa de labor}
\textsuperscript{(38)} o \VUgras{dedal,} o \VUgras{fío} \textsuperscript{(36)}, as
\VUgras{agullas} \textsuperscript{(37)} e as \VUgras{tesoiras}
\textsuperscript{(39)}, e pechar coa \VUgras{chave} \textsuperscript{(40)} a tapa da
mesa.

%%K t06-c1-09-1 p
9. --- Sobre o \VUgras{veladoiro} \textsuperscript{(41)} do medio, Xustina renovará a
auga da \VUgras{garrafa} \textsuperscript{(42)}. Porá \VUgras{azucre} no
\VUgras{azucreiro} \textsuperscript{(43)}, limpará a \VUgras{pinza do azucre}
\textsuperscript{(44)} e retirará o \VUgras{cepillo da roupa} \textsuperscript{(46)}.

%%K t06-c1-10-1 p
10. --- O lado dereito do cuarto pediralle menos traballo, pois a \VUgras{cadeira
longa} \textsuperscript{(47)} e o seu \VUgras{coxín} \textsuperscript{(48)} están no
seu sitio xusto e, como non fai frío, nada está fóra de lugar entre os apeiros da
\VUgras{lareira} \textsuperscript{(49)}. A \VUgras{pa} \textsuperscript{(50)}, as
\VUgras{tenaces} \textsuperscript{(51)}, os \VUgras{morillos}
\textsuperscript{(55)} e o \VUgras{gardacinzas} \textsuperscript{(56)} están limpos;
a \VUgras{pantalla do lume} \textsuperscript{(54)} non foi movida e a \VUgras{caixa
da leña} \textsuperscript{(52)} está ben provista de \VUgras{leña}
\textsuperscript{(53)}.

%%K t06-noto-1 noto
\VUnotes{143.2mm}{%
(*) Babo = neno pequeno; Ing. \textit{baby}, Fr. \textit{bébé}, It.
\textit{bambino}.%
}

%%K t06-noto-2 noto
\VUnotes{143.2mm}{%
(*) Lambrequino = Fr. \textit{lambrequin}, Esp. \textit{lambrequines}, Port.
\textit{lambrequins}, e mesmo Al. e Ing. \textit{lambrequins}; polo tanto en alemán,
inglés, francés, portugués e español.%
}

%%K t06-c1-11-1 p
11. --- Aínda así, terá que quitar o po ao \VUgras{reloxo de repisa}
\textsuperscript{(57)}, aos \VUgras{candelabros} \textsuperscript{(58)} e aos
\VUgras{baleirapetos} \textsuperscript{(59)}, e finalmente limpar o
\VUgras{espello} \textsuperscript{(60)}.

%%K t06-c1-12-1 p
12. --- Canto ao \VUgras{berce} \textsuperscript{(61)} do neniño (do bebé (*)), ese
xa está listo.

%%K t06-tit-3 sub
\VUsaut{5.0mm}
\VUpk{10.2pt}{40}{O salón.}
\VUsaut{3.0mm}

%%K t06-c2-01-1 p
1. --- Velaquí agora o salón. É un \VUgras{cuarto} moi belo. A lareira, que está no
recanto dereito, vai adornada cunha delicada \VUgras{estatuíña}
\textsuperscript{(1)} e con dous belos \VUgras{floreiros} \textsuperscript{(3)}, e
vestida cun \VUgras{lambrequín} de veludo \textsuperscript{(2)} (*).

%%K t06-c2-02-1 p
2. --- Para impedir que as faíscas do lar prendan na alfombra, puxeron diante do lar
un \VUgras{gardafaíscas} \textsuperscript{(4)} de cobre.

%%K t06-c2-03-1 p
3. --- Preto de alí está aberto un elegante \VUgras{escritorio} de señora
\textsuperscript{(5)}. Nel a \VUgras{dona da casa} \textsuperscript{(23)} escribe as
súas cartas.

%%K t06-c2-04-1 p
4. --- A fiestra vai adornada con \VUgras{cortinas de vidro} \textsuperscript{(6)}
con \VUgras{alzapanos} \textsuperscript{(8)} enganchados aos \VUgras{ganchos}
\textsuperscript{(7)}, e con cortinóns de tea \VUgras{recollidos por riba}
\textsuperscript{(9)}.

%%K t06-c2-05-1 p
5. --- No van da fiestra, unha \VUgras{floreira} \textsuperscript{(11)} garda unha
\VUgras{planta} verde \textsuperscript{(10)}, unha soberbia palmeira cuxas follas se
estenden case ata a \VUgras{lámpada de petróleo} \textsuperscript{(12)}.

%%K t06-c2-06-1 p
6. --- A \VUgras{pantalla} \textsuperscript{(13)} da lámpada foi arranxada por María,
a encantadora \VUgras{moza} \textsuperscript{(14)} que está ao piano
\textsuperscript{(15)}. Sentada moi dereita no \VUgras{tallo}
\textsuperscript{(16)}, estuda unha peza de música. É moi hábil e coidadosa, como o
proba o seu \VUgras{armario da música} \textsuperscript{(17)}, que está
perfectamente ordenado.

%%K t06-tit-4 sub
\VUsaut{5.0mm}
\VUpk{10.2pt}{40}{O rapaz travieso.}
\VUsaut{3.0mm}

%%K t06-c2-07-1 p
7. --- O seu irmán, Paulo, busca un \VUgras{libro} \textsuperscript{(19)} na
\VUgras{biblioteca} \textsuperscript{(18)}. É un \VUgras{mozo}
\textsuperscript{(20)} inquedo e insoportable. Ao agarrarse á biblioteca para
acadar o libro, que está alto de máis, arrisca facer caer o \VUgras{busto}
\textsuperscript{(21)} que está enriba.

%%K t06-c2-08-1 p
8. --- Aproveita, para facer esa parvada, que a súa nai está ocupada conversando cunha
amiga, unha \VUgras{dama} \textsuperscript{(22)} que veu pasar uns días na súa casa. A
nai senta nun \VUgras{sofá} \textsuperscript{(24)} preto do \VUgras{sillón}
\textsuperscript{(25)}, e a súa amiga está no \VUgras{asento}
\textsuperscript{(27)}, do que só se ve o \VUgras{respaldo} \textsuperscript{(26)}.

%%K t06-c2-09-1 p
9. --- O gran \VUgras{marco} \textsuperscript{(30)} que pendura na parede garda o
\VUgras{retrato} \textsuperscript{(29)} dunha parenta. No \VUgras{álbum}
\textsuperscript{(32)} pousado na mesa están xuntadas as \VUgras{fotografías}
\textsuperscript{(33)} dos demais membros da familia. Os nenos acaban de o follear,
pois as \VUgras{cadeiras} \textsuperscript{(31)} seguen agrupadas arredor da mesa.

%%K t06-c2-10-1 p
10. --- O \VUgras{floreiro chinés} \textsuperscript{(34)} de porcelana, que está preto
do \VUgras{abrecartas} \textsuperscript{(35)}, foille regalado a María polo seu
santo, e o \VUgras{biombo} \textsuperscript{(36)} que adorna o recanto do cuarto
tróuxoo da China o seu tío.

%%K t06-c2-11-1 p
11. --- Alegrado polos fermosísimos \VUgras{cadros} \textsuperscript{(37)}
pendurados no \VUgras{tabique} \textsuperscript{(42)}, alumado pola \VUgras{lámpada
do teito} \textsuperscript{(38)}, cuxas \VUgras{lámpadas eléctricas}
\textsuperscript{(39)} brillan ledas á noite, ese salón parece moi agradable. A branda
\VUgras{alfombra} \textsuperscript{(40)} estendida no parqué, e o \VUgras{timbre
eléctrico} \textsuperscript{(41)} tan doado de usar, rematan o seu confort.

%%K t06-tit-5 sub
\VUsaut{3.6mm}
\VUpk{10.2pt}{40}{O comedor.}
\VUsaut{3.0mm}

%%K t06-c3-01-1 p
1. --- Tocou a campá da cea. Toda a familia, agás María, está xuntada arredor da mesa.
O comedor está agradablemente quente grazas á excelente \VUgras{estufa}
\textsuperscript{(1)} de louza, co seu \VUgras{tubo} delgado \textsuperscript{(2)}.

%%K t06-c3-02-1 p
2. --- Ese cuarto non ten moitos mobles. Ao longo da parede pendura, por riba da
\VUgras{repisa} \textsuperscript{(4)}, unha \VUgras{fonteciña}
\textsuperscript{(3)} de adorno. Moi preto está o \VUgras{aparador}
\textsuperscript{(5)}, no que se pon a comida fría, os pratos da sobremesa e as
distintas pezas do servizo que non fan falta de contado.

%%K t06-c3-03-1 p
3. --- Así vemos no estante de arriba un \VUgras{polo asado} \textsuperscript{(7)},
\VUgras{peixe} \textsuperscript{(8)} e unha \VUgras{copa} \textsuperscript{(10)} con
\VUgras{froitas} \textsuperscript{(9)}. Debaixo velaquí o \VUgras{queixo}
\textsuperscript{(11)}, o \VUgras{pan} \textsuperscript{(12)}, as
\VUgras{vinagreiras} \textsuperscript{(13)} que teñen todo o que fai falta para
temperar a ensalada: o aceite, o vinagre, o \VUgras{sal} \textsuperscript{(14)} e a
\VUgras{pementa} \textsuperscript{(15)}. Finalmente, no estante de abaixo hai un
\VUgras{bolo} \textsuperscript{(17)} preto da \VUgras{mostardeira}
\textsuperscript{(16)}.

%%K t06-c3-04-1 p
4. --- O reloxo do comedor, ao que se chama \VUgras{reloxo de parede}
\textsuperscript{(20)}, ten unha forma moi especial. Vai encaixado nun marco de
madeira que garda ademais un \VUgras{barómetro} \textsuperscript{(19)} e un
\VUgras{termómetro} \textsuperscript{(18)}.

%%K t06-tit-6 sub
\VUsaut{4.6mm}
\VUpk{10.2pt}{40}{Como se serve a cea.}
\VUsaut{3.0mm}

%%K t06-c3-05-1 p
5. --- En todas as paredes do cuarto están postos \VUgras{paneis de madeira}
\textsuperscript{(21)} de carballo encerado. Unha \VUgras{cortina de porta}
\textsuperscript{(22)} de tea pecha abondo a porta que dá entrada ao alpendre. Alí a
familia irá tomar o café despois da cea. Xa as \VUgras{cuncas}
\textsuperscript{(24)} e os \VUgras{pratiños} \textsuperscript{(25)} están listos no
\VUgras{moble da louza} \textsuperscript{(23)}.

%%K t06-c3-06-1 p
6. --- Por riba da mesa vai fixada ao teito unha \VUgras{lámpada colgante}
\textsuperscript{(26)} cuxa luz moi viva alumea á noite todo o comedor.

%%K t06-c3-07-1 p
7. --- Vexamos agora a \VUgras{mesa de comer} \textsuperscript{(27)} mesma. Cando
entrou a familia, o servizo estaba posto. Sobre o \VUgras{mantel}
\textsuperscript{(28)} estaban postos, no sitio de cada comensal, un \VUgras{prato}
\textsuperscript{(33)}, un \VUgras{gardanapo} \textsuperscript{(29)}, unha
\VUgras{culler} \textsuperscript{(34)}, un \VUgras{garfo} \textsuperscript{(35)}, un
\VUgras{coitelo} \textsuperscript{(36)} e un \VUgras{vaso} \textsuperscript{(32)}.
Había tamén \VUgras{botellas} \textsuperscript{(30)} para o viño e unha
\VUgras{garrafa} \textsuperscript{(31)} para a auga.

%%K t06-c3-08-1 p
8. --- O \VUgras{criado} \textsuperscript{(39)} trouxo primeiro a \VUgras{sopeira}
\textsuperscript{(37)}, na que aínda fumega algo de sopa. Agora serve o primeiro
\VUgras{prato} \textsuperscript{(38)}, un prato de carne. Despois presentará os que xa
nomeamos e, finalmente, despois do \VUgras{prato da sobremesa}, aproveitará que os
seus donos entrasen no alpendre para retirar os apeiros da mesa e volver arranxar o
comedor.

%%K t06-c3-08-2 p
\textit{Nota.} --- \textit{As palabras correntes que se refiren aos alimentos
indícanse aquí moi en resumo. A lista completarase nos dous cadros «O hotel» (n.º 13)
e «O mercado» (n.º 15).}

%%K t06-tit-7 sub
\VUsaut{4.6mm}
\VUpk{10.2pt}{40}{A cociña.}
\VUsaut{3.0mm}

%%K t06-c4-01-1 p
1. --- Na cociña, que ollamos pola porta entreaberta, vemos unha desorde pintoresca.
Diante do \VUgras{lar} \textsuperscript{(1)}, no que un \VUgras{lume}
\textsuperscript{(3)} crepita, está instalado un \VUgras{asador xiratorio}
\textsuperscript{(5)}. Un fermoso \VUgras{asado} \textsuperscript{(7)} está
\VUgras{espetado} \textsuperscript{(6)} e vai cocendo.

%%K t06-c4-02-1 p
2. --- O \VUgras{fol} \textsuperscript{(8)} e a \VUgras{pa do carbón}
\textsuperscript{(9)}, que acaban de usar, quedaron no chan, o mesmo ca a
\VUgras{leña} \textsuperscript{(10)}.

%%K t06-c4-03-1 p
3. --- O alto da lareira está ordenado; a \VUgras{lanterna} \textsuperscript{(11)}, o
\VUgras{misteiro} \textsuperscript{(13)}, no que están os \VUgras{mistos}
\textsuperscript{(12)}, o \VUgras{candeeiro} \textsuperscript{(14)} e a
\VUgras{palmatoria} \textsuperscript{(15)} coa súa \VUgras{vela}
\textsuperscript{(16)} están no seu sitio. Pero o gran \VUgras{balde do carbón}
\textsuperscript{(18)}, cheo do \VUgras{carbón de pedra} \textsuperscript{(17)} que a
\VUgras{cociñeira} \textsuperscript{(23)} acaba de encher no soto, ficou no medio da
cociña.

%%K t06-c4-04-1 p
4. --- Abofé que a pobre muller ten que apurar para que o xantar estea listo a tempo.
Agora está preto da \VUgras{cociña económica} \textsuperscript{(19)} e bate unha
\VUgras{salsa} \textsuperscript{(24)} que non se debe queimar de máis.

%%K t06-c4-05-1 p
5. --- No \VUgras{forno} \textsuperscript{(20)} coce un bolo que preparou para a
merenda dos nenos. Por riba da cociña penduran a \VUgras{culler grande}
\textsuperscript{(21)} e a \VUgras{escumadeira} \textsuperscript{(22)}.

%%K t06-tit-8 sub
\VUsaut{4.0mm}
\VUpk{10.2pt}{40}{A louza.}
\VUsaut{3.0mm}

%%K t06-c4-06-1 p
6. --- A \VUgras{batería de cociña} \textsuperscript{(25)}, colgada ao fondo do
cuarto, está moi limpa. As belas \VUgras{cazolas} de cobre \textsuperscript{(26)}, a
\VUgras{leiteira} \textsuperscript{(27)}, o \VUgras{coador} \textsuperscript{(28)} e
o \VUgras{relador} \textsuperscript{(29)} foron limpados con coidado.

%%K t06-c4-07-1 p
7. --- A \VUgras{caixa do sal} \textsuperscript{(30)} está posta no moble da louza
preto da \VUgras{teteira} \textsuperscript{(31)} e da \VUgras{botella grande do
aceite} \textsuperscript{(32)}. O \VUgras{caixón} \textsuperscript{(33)} garda de
seguro os cubertos e os coitelos de cociña.

%%K t06-c4-08-1 p
8. --- A \VUgras{vasoira} \textsuperscript{(34)} e o \VUgras{plumeiro}
\textsuperscript{(35)} están nun recanto. A cociñeira acaba de usar o \VUgras{bloque
de madeira} \textsuperscript{(36)}; deixouno preto da fiestra.

%%K t06-c4-09-1 p
9. --- Unha \VUgras{toalla de mans} \textsuperscript{(37)} pendura na parede, preto do
\VUgras{funil} \textsuperscript{(38)}. Nesa parte da cociña están gardados a
\VUgras{grella} \textsuperscript{(39)}, o \VUgras{picador}
\textsuperscript{(40)} e a \VUgras{táboa de picar} \textsuperscript{(41)}. Despois,
por riba do picador, nun \VUgras{estante} \textsuperscript{(42)}, hai
\VUgras{potes} \textsuperscript{(43)}, a \VUgras{pota de ferver}
\textsuperscript{(44)} coa súa \VUgras{tapa} \textsuperscript{(45)} e a
\VUgras{marmita} \textsuperscript{(46)}. A \VUgras{tixola} \textsuperscript{(47)}
pendura preto.

%%K t06-c4-10-1 p
10. --- Só a \VUgras{billa} \textsuperscript{(48)} de cobre bota brillo neste recanto.
O \VUgras{fregadeiro} \textsuperscript{(49)}, no que está pousada a grosa
\VUgras{xerra} \textsuperscript{(50)}, debe de ser moi cómodo co seu armario
pequeno, no que se garda o \VUgras{barril do vinagre} \textsuperscript{(51)} provisto
da súa \VUgras{billa} \textsuperscript{(52)}, a \VUgras{caixa dos restos}
\textsuperscript{(53)} e os \VUgras{pratos sucios} \textsuperscript{(54)}.

%%K t06-tit-9 sub
\VUsaut{4.0mm}
\VUpk{10.2pt}{40}{Outros obxectos que hai na cociña.}
\VUsaut{3.0mm}

%%K t06-c4-11-1 p
11. --- A \VUgras{cuba} \textsuperscript{(55)} na que se lavan os \VUgras{legumes}
\textsuperscript{(58)} e a \VUgras{caldeira} de ferro fundido \textsuperscript{(56)}
pousada no \VUgras{chan de baldosas} \textsuperscript{(57)} deben de ter tamén o seu
sitio nese armario.

%%K t06-c4-12-1 p
12. --- De certo que á cociñeira non lle faltan fogóns, pois velaquí ademais, no
recanto da mesa, unha \VUgras{cociña de gas} \textsuperscript{(59)} na que se quenta
auga nunha cazola pequena. O \VUgras{vapor} \textsuperscript{(60)} que sae dela
indícanos que debe de estar a ferver. Seguramente esa auga se usará para o café, pois
velaquí no outro extremo da mesa a \VUgras{cafeteira} \textsuperscript{(64)} e o
\VUgras{muíño do café} \textsuperscript{(63)}.

%%K t06-c4-13-1 p
13. --- A cociña de gas vai unida ao \VUgras{queimador de gas}
\textsuperscript{(62)} por un longo \VUgras{tubo de goma} \textsuperscript{(61)}.

%%K t06-c4-14-1 p
14. --- A \VUgras{criada de día} \textsuperscript{(66)} que vén axudar a cociñeira
prepara os legumes para a cea. Cando estean limpados e lavados, poñeraos na
\VUgras{bacía} \textsuperscript{(65)} agardando a hora de os cocer.

%%K t06-tit-10 sub
\VUsaut{4.0mm}
{\centering\fontsize{10.2pt}{10.2pt}\selectfont\textls[40]{\scshape O cuarto de baño.} \textit{(Cuarto no que un se baña.)}\par}
\VUsaut{3.0mm}

%%K t06-c5-01-1 p
1. --- Só nos queda unha indiscreción por facer para coñecermos os cuartos principais
da vivenda: ver a instalación do cuarto de baño. Iso non levará moito tempo, pois non
é grande, e saberemos axiña que a \VUgras{bañeira} \textsuperscript{(1)} ten un bo
\VUgras{quentador} \textsuperscript{(2)}, que a \VUgras{xaboeira}
\textsuperscript{(3)} está ao alcance da man de quen se baña, e que o
\VUgras{toalleiro} \textsuperscript{(5)} ha facer doado o secado das
\VUgras{toallas} \textsuperscript{(4)}.

%%K t06-c5-02-1 p
2. --- Ese cuarto ten as paredes cubertas de \VUgras{azulexos}
\textsuperscript{(6)}, o que permitiu instalar unha \VUgras{ducha}
\textsuperscript{(7)} sen medo ningún de que a auga, que salpica mentres se usa,
estrague os muros. Unha \VUgras{baciña} de cinc \textsuperscript{(8)}, que serve para
os baños de pés, completa o mobiliario.
\VUsaut{6.0mm}
\VUfilet{22mm}
